\subsection{Dataset Description and Experimental Setup}

The empirical foundation of this thesis rests on a real-world smart manufacturing sensor dataset comprising 100,000 timestamped observations collected from a group of 50 industrial machines over a continuous 69-day monitoring period spanning from January 1, 2025 at 00:00:00 to March 11, 2025 at 10:39:00. Each observation records the simultaneous state of thirteen variables: a timestamp, a machine identifier (ranging from 1 to 50), five continuous sensor channels (temperature, vibration, humidity, pressure, and energy consumption), three derived condition indicators (machine status, anomaly flag, and downtime risk), and three maintenance-relevant labels (predicted remaining useful life, failure type, and maintenance required). The dataset is complete with no missing values across all 1,300,000 individual data points, a property that reflects the robust data quality management practices of the industrial monitoring system from which it was collected.

The dataset directly operationalizes three distinct experimental pipelines implemented across the three QASAMAP code components. The full 100,000-observation dataset is loaded, feature-engineered, and partitioned into per-machine training (80\%) and test (20\%) splits for T-GCN and comparative model training and evaluation. A representative five-event sample from Machine 40 spanning 10:40:00 to 10:44:00 on March 11, 2025 serves as the streaming inference demonstration. A multi-machine sample of 11 events covering machines 1, 2, 3, 4, 32, 40, and 47 provides the basis for the full quantum-agentic pipeline demonstration and three-way comparative analysis. All experiments were executed on a server equipped with an NVIDIA A100-SXM4-80GB GPU, which provided the computational environment for deep learning model training in the forecasting and classification experiments.

The experimental evaluation is organized into four distinct components: (1) spatiotemporal forecasting model comparison across T-GCN and LSTM-based architectures; (2) machine health classification and regression across five supervised learning algorithms; (3) agentic AI pipeline evaluation without quantum enhancement; and (4) agentic AI pipeline evaluation with full quantum computing integration. Together, these four components provide a complete empirical characterisation of the QASAMAP framework, enabling both component-level analysis and full-system performance attribution.

\subsection{Sensor and Target Variable Summary}

The five continuous sensor channels exhibit the statistical profiles summarized in Table~\ref{tab:sensor_stats}. Temperature is the most discriminative sensor for anomaly detection, with a mean of 75.02\textdegree C and a Pearson correlation of $r = 0.447$ with the anomaly flag. Vibration shows a 19\% elevation in anomalous readings (mean 58.51) compared to normal readings (mean 49.18) and is the primary indicator of the Vibration Issue failure type. Humidity exhibits near-zero linear correlation with anomaly occurrence, contributing to fault detection through nonlinear interaction effects captured by the deep learning layers. Pressure and energy consumption similarly show weak linear correlations with anomaly occurrence but contribute to the detection of Pressure Drop and Electrical Fault failure types respectively.

\begin{table}[htbp]
\centering
\caption{Descriptive Statistics of Sensor Variables ($n = 100{,}000$)}
\label{tab:sensor_stats}
\begin{tabular}{lccccc}
\hline
\textbf{Sensor} & \textbf{Min} & \textbf{Max} & \textbf{Mean} & \textbf{Std Dev} & \textbf{Corr.\ w/ Anomaly} \\
\hline
Temperature (\textdegree C) & 35.55 & 121.94 & 75.02 & 10.03 & 0.447 \\
Vibration         & $-$17.09 & 113.80 & 50.42  & 14.31 & 0.198 \\
Humidity (\%)     & 30.00  & 80.00  & 55.01  & 10.00 & $\approx 0.001$ \\
Pressure (bar)    & 1.00   & 5.00   & 3.00   & 1.15  & $\approx -0.002$ \\
Energy (kWh)      & 0.50   & 5.00   & 2.75   & 1.30  & $\approx 0.006$ \\
\hline
\end{tabular}
\end{table}

The target variable distributions are summarized in Table~\ref{tab:target_stats}. The binary anomaly flag shows an 8.92\% positive rate, which is characteristic of real industrial anomaly detection problems and drives the choice of F1-score as the primary evaluation metric. The failure type variable exhibits severe class imbalance, with Normal operation constituting 91.90\% of readings. The predicted remaining useful life (RUL) variable follows a near-uniform distribution across the range 1--499 hours, with 16,333 critical readings below 50 hours and 1,637 readings below 10 hours representing machines on the verge of imminent failure.

\begin{table}[htbp]
\centering
\caption{Target Variable Distributions}
\label{tab:target_stats}
\begin{tabular}{lcc}
\hline
\textbf{Target Variable} & \textbf{Value / Class} & \textbf{Count (\%)} \\
\hline
Anomaly Flag         & Anomalous (1)            & 8,916 (8.92\%)  \\
                     & Normal (0)               & 91,084 (91.08\%) \\
\hline
Machine Status       & Running (1)              & 80,091 (80.09\%) \\
                     & Warning (2)              & 10,057 (10.06\%) \\
                     & Stopped (0)              & 9,852 (9.85\%)  \\
\hline
Failure Type         & Normal                   & 91,899 (91.90\%) \\
                     & Vibration Issue          & 3,129 (3.13\%)  \\
                     & Overheating              & 1,989 (1.99\%)  \\
                     & Pressure Drop            & 1,969 (1.97\%)  \\
                     & Electrical Fault         & 1,014 (1.01\%)  \\
\hline
Downtime Risk        & High Risk ($= 1.0$)      & 8,904 (8.90\%)  \\
                     & No Risk ($= 0.0$)        & 91,084 (91.08\%) \\
\hline
Maintenance Required & Required (1)             & 19,700 (19.70\%) \\
                     & Not Required (0)         & 80,300 (80.30\%) \\
\hline
RUL (hours)          & Mean: 234.27 h           & Std: 150.06 h   \\
                     & Critical ($< 50$ h)      & 16,333 (16.33\%) \\
                     & Imminent ($< 10$ h)      & 1,637 (1.64\%)  \\
\hline
\end{tabular}
\end{table}

\subsection{Forecasting Model Experiment}

The first experimental component evaluates spatiotemporal sensor forecasting performance across three architectures: a standalone Optimized LSTM baseline, the proposed T-GCN with cross-correlation adjacency (CC-Adj), and a Hybrid LSTM-GCN combining both architectures. All models were trained and evaluated on the 100,000-observation dataset partitioned into 80/20 per-machine train/test splits. The cross-correlation graph for T-GCN comprises 50 nodes (machines) and 538 directed edges constructed using the \texttt{compute\_machine\_lag\_matrix} function. Machine 40 was identified as the anchor node based on its highest out-degree centrality, with key top influencer machines being 26, 38, 42, 47, and 14, and key top dependent machines being 29, 40, 28, 8, and 16. Training was performed on an NVIDIA A100-SXM4-80GB GPU.

The per-feature forecasting results across the five sensor channels and all three models are presented in Table~\ref{tab:forecasting_results}. Summary performance metrics and training times are presented in Table~\ref{tab:forecasting_summary}.

\begin{table}[htbp]
\centering
\caption{Per-Feature Forecasting Results: MSE, RMSE, MAE, and R\textsuperscript{2} for Each Model and Sensor}
\label{tab:forecasting_results}
\resizebox{\textwidth}{!}{%
\begin{tabular}{llcccc}
\hline
\textbf{Model} & \textbf{Feature} & \textbf{MSE} & \textbf{RMSE} & \textbf{MAE} & \textbf{R\textsuperscript{2}} \\
\hline
\multirow{5}{*}{OptimizedLSTM}
  & temperature      & 3,035.27  & 55.09 & 45.42 & 0.6993 \\
  & vibration        & 8,618.87  & 92.84 & 76.38 & 0.5735 \\
  & humidity         & 1,151.30  & 33.93 & 27.81 & 0.8761 \\
  & pressure\_sensor & 34,765.57 & 186.45& 153.13& 0.4254 \\
  & energy\_consumption & 19,691.01 & 140.32 & 117.43 & 0.5522 \\
\hline
\multirow{5}{*}{T-GCN (CC-Adj)}
  & temperature      & 2,658.46  & 51.56 & 42.00 & 0.7366 \\
  & vibration        & 8,200.14  & 90.55 & 73.99 & 0.5942 \\
  & humidity         & 855.65    & 29.25 & 23.83 & 0.9080 \\
  & pressure\_sensor & 30,286.70 & 174.03& 142.04& 0.5002 \\
  & energy\_consumption & 15,339.71 & 123.85 & 101.04 & 0.6581 \\
\hline
\multirow{5}{*}{Hybrid LSTM-GCN}
  & temperature      & 3,127.64  & 55.93 & 45.93 & 0.6901 \\
  & vibration        & 9,049.22  & 95.13 & 78.51 & 0.5522 \\
  & humidity         & 989.76    & 31.46 & 25.72 & 0.8938 \\
  & pressure\_sensor & 33,780.66 & 183.80& 150.84& 0.4417 \\
  & energy\_consumption & 21,482.71 & 146.57 & 122.29 & 0.4894 \\
\hline
\end{tabular}%
}
\end{table}

\begin{table}[htbp]
\centering
\caption{Forecasting Model Summary: Average Metrics and Training Time}
\label{tab:forecasting_summary}
\begin{tabular}{lcccccc}
\hline
\textbf{Model} & \textbf{Avg MSE} & \textbf{Avg RMSE} & \textbf{Avg MAE} & \textbf{Avg R\textsuperscript{2}} & \textbf{Train Time (s)} \\
\hline
OptimizedLSTM     & 13,452.40 & 101.73 & 84.03 & 0.6853 & 52.56 \\
T-GCN (CC-Adj)    & 11,468.13 & 93.85  & 76.58 & 0.7702 & 111.28 \\
Hybrid LSTM-GCN   & 13,365.80 & 102.58 & 84.66 & 0.7134 & 47.89 \\
\hline
\end{tabular}
\end{table}

The T-GCN with cross-correlation adjacency achieves the highest average R\textsuperscript{2} of 0.7702, representing a 12.4\% improvement over the OptimizedLSTM baseline (R\textsuperscript{2} = 0.6853) and an 8.0\% improvement over the Hybrid LSTM-GCN (R\textsuperscript{2} = 0.7134). The improvement is most pronounced for humidity prediction (R\textsuperscript{2} = 0.9080 vs.\ 0.8761 for LSTM) and energy consumption (R\textsuperscript{2} = 0.6581 vs.\ 0.5522), where the graph-based propagation of correlated cross-machine signals provides the greatest benefit. The pressure sensor channel remains the most challenging to forecast across all three models, with R\textsuperscript{2} values ranging from 0.4254 to 0.5002, reflecting the lower unique-value count (401 distinct values) and the more complex, possibly setpoint-driven dynamics of pressure in the monitored machine pool.

Although the T-GCN requires a longer training time (111.28 seconds vs.\ 52.56 seconds for the standalone LSTM), this overhead is attributable to the graph convolution operations over 538 edges during each forward pass and is acceptable given the significant performance gains. The Hybrid LSTM-GCN achieves intermediate R\textsuperscript{2} performance (0.7134) with the shortest training time (47.89 seconds), offering a practical trade-off when computational resources are constrained.

\subsubsection{Forecasting Experiment Visualisations}

\paragraph{Cross-Correlation Strength Matrix.}
Figure~\ref{fig:cc_matrix} presents the pairwise cross-correlation strength matrix computed from the temperature time series of the first 20 machines in the dataset. Each cell $(i, j)$ encodes the absolute Pearson correlation between machine $i$ and machine $j$ at the lag of maximum correlation. The diagonal entries are identically 1.0 (each machine is perfectly correlated with itself), while off-diagonal entries reveal the inter-machine temperature coupling structure that defines the T-GCN adjacency graph. The predominance of light-yellow and light-orange cells indicates that most machine pairs exhibit weak to negligible direct linear coupling ($|r| < 0.35$), consistent with independent production lines operating under shared environmental conditions. Scattered deeper-orange and red cells identify specific machine pairs --- notably machines 3--13, 17--19, and 19--20 --- with moderate negative correlations, suggesting thermal counter-cycling patterns in which one machine's heat signature inversely predicts another's at a characteristic lag. Only machine pairs with $|r| > 0.35$ are retained as directed edges in the T-GCN graph, yielding the 538 edges that constitute the spatial backbone of the model.

\begin{figure}[htbp]
    \centering
    \includegraphics[width=0.85\textwidth]{figures/cross-correlation-strength-matrix.png}
    \caption{Cross-Correlation Strength Matrix for temperature time series across 20 machines. Each cell represents the Pearson correlation at the lag of maximum absolute correlation. Diagonal entries are 1.0 (self-correlation). Only pairs with $|r| > 0.35$ are retained as directed edges in the T-GCN cross-correlation adjacency graph. Predominantly weak off-diagonal correlations confirm that most machines operate independently, with isolated pairs exhibiting moderate negative coupling that T-GCN exploits for predictive propagation.}
    \label{fig:cc_matrix}
\end{figure}

\paragraph{Model Performance Comparison.}
Figure~\ref{fig:model_performance} presents a four-panel performance summary comparing the three forecasting architectures across R\textsuperscript{2} per feature, MSE per feature, aggregate metrics, and training loss convergence. The top-left panel confirms T-GCN's consistent R\textsuperscript{2} advantage across all seven sensor channels, with the most pronounced improvements observed for vibration (R\textsuperscript{2}: 0.594 vs.\ 0.574 for LSTM) and energy consumption (R\textsuperscript{2}: 0.658 vs.\ 0.552 for LSTM). The engineered features \texttt{temp\_vibe\_id} and \texttt{press\_energy} show competitive R\textsuperscript{2} scores near 0.7--0.8 across all models, confirming that domain-driven feature construction provides useful complementary predictive signals. The top-right MSE panel highlights that the \texttt{temp\_vibe\_id} feature dominates absolute error magnitude due to its large numerical scale (product of temperature and vibration values), while all other sensor channels produce near-zero MSE on the visible scale. The bottom-left aggregate metrics panel clearly shows T-GCN achieving the highest average R\textsuperscript{2} ($\approx$0.77), with the Hybrid LSTM-GCN intermediate ($\approx$0.71) and the standalone LSTM lowest ($\approx$0.69). The bottom-right training loss convergence panel is particularly informative: T-GCN converges most rapidly, dropping steeply between epochs 0 and 4 and stabilising at a final loss of $\approx$0.024, substantially below the OptimizedLSTM final loss of $\approx$0.031 and the Hybrid LSTM-GCN loss of $\approx$0.028. This rapid convergence reflects the graph convolution layers' ability to immediately leverage cross-machine spatial structure from epoch 1, in contrast to the purely temporal LSTM that must build up temporal representations more gradually.

\begin{figure}[htbp]
    \centering
    \includegraphics[width=\textwidth]{figures/model-performance-lstm-tgcn-hybrid.png}
    \caption{Four-panel model performance comparison: (top-left) R\textsuperscript{2} score per sensor feature for all three architectures; (top-right) MSE per feature on log scale; (bottom-left) aggregate average R\textsuperscript{2}, inverse MSE, and inverse MAE; (bottom-right) training loss convergence over 15 epochs. T-GCN (CC-Adj, gold) consistently outperforms both the OptimizedLSTM (teal) and Hybrid LSTM-GCN (magenta) baselines and achieves the fastest convergence to the lowest final training loss.}
    \label{fig:model_performance}
\end{figure}

\paragraph{Multi-Horizon Autoregressive Prediction.}
Figure~\ref{fig:multi_horizon} presents a $4 \times 3$ panel grid showing the autoregressive rollout prediction performance across 10 forecast horizon steps for four sensor channels (temperature, vibration, humidity, pressure) and all three models. Each subplot overlays the actual sensor trajectory (solid black line), the model's predicted trajectory (dashed coloured line), and a shaded $\pm 1\sigma$ uncertainty band computed across multiple starting windows. Several observations are noteworthy. First, all three models track the actual temperature and vibration trajectories closely at short horizons (steps 1--3), with divergence increasing beyond step 5 as autoregressive error accumulates. Second, the T-GCN uncertainty bands (gold shading) are generally narrower than those of the OptimizedLSTM (blue shading) for the humidity and pressure channels, indicating more consistent multi-step predictions across different starting windows. Third, the humidity channel shows the strongest overall predictive accuracy across all models, consistent with its highest individual R\textsuperscript{2} scores in the per-feature benchmark. The pressure channel exhibits the largest $\pm 1\sigma$ spread, reflecting the setpoint-driven dynamics that resist smooth autoregressive rollout by all architectures.

\begin{figure}[htbp]
    \centering
    \includegraphics[width=\textwidth]{figures/multi-horizon-per-model-feature.png}
    \caption{Multi-horizon autoregressive rollout predictions for four sensor channels (rows: temperature, vibration, humidity, pressure) across three models (columns: OptimizedLSTM, T-GCN CC-Adj, Hybrid LSTM-GCN). Solid black line: actual values; dashed coloured line: predicted values; shaded band: $\pm 1\sigma$ uncertainty across starting windows. T-GCN consistently maintains tighter uncertainty bands and closer tracking of actual trajectories, particularly for humidity and pressure, demonstrating superior multi-step robustness compared to baseline architectures.}
    \label{fig:multi_horizon}
\end{figure}

\paragraph{Horizon MAE Degradation Curve.}
Figure~\ref{fig:mae_degradation} quantifies forecast error accumulation as a function of prediction horizon. The four subplots present MAE vs.\ horizon step (1--10) for temperature, vibration, humidity, and pressure respectively. A clear pattern of non-monotonic MAE growth is observed across all channels: MAE does not increase uniformly with horizon but instead exhibits local peaks at steps 4--6, followed by partial recovery at longer horizons. This non-monotonic pattern reflects the semi-periodic nature of the sensor dynamics, where the model's autoregressive predictions re-synchronise with the actual trajectory as it returns to its mean. T-GCN consistently achieves the lowest MAE at almost every horizon step across all four channels, with the margin over the LSTM baseline widening at the critical medium-range horizons (steps 5--8) that are most operationally relevant for predictive maintenance scheduling. The vibration channel shows the highest absolute MAE values (ranging from approximately 1.2 to 2.0 across models), confirming it as the most difficult channel for multi-step forecasting, while humidity achieves the lowest MAE despite its wide dynamic range.

\begin{figure}[htbp]
    \centering
    \includegraphics[width=\textwidth]{figures/horizon-mae-degradation-curve.png}
    \caption{Horizon MAE Degradation Curves showing forecast error accumulation over 10 prediction steps for temperature, vibration, humidity, and pressure channels. All three models exhibit non-monotonic MAE growth reflecting sensor periodicity. T-GCN (CC-Adj, gold) consistently achieves the lowest MAE at most horizon steps, with the advantage most pronounced at medium horizons (steps 5--8) that are most critical for maintenance decision scheduling.}
    \label{fig:mae_degradation}
\end{figure}

\paragraph{Multi-Source Cross-Correlation Graph Simulation.}
Figure~\ref{fig:cc_graph_sim} illustrates the T-GCN's cross-correlation graph propagation mechanism through a multi-source simulation using three source machines (Machine 40, Machine 27, and Machine 14) and three dependent machines (Machine 1, Machine 50, and Machine 32). During the first 50 time steps, the three source machines (dashed lines) exhibit high-amplitude oscillations between approximately 50 and 100 units, representing normal operating variability in the simulated sensor feature. The dependent machines (solid lines) mirror these oscillations with a characteristic lag, visually demonstrating the propagation of cross-machine influence through the graph edges. At approximately time step 50, all traces collapse to near-zero, simulating the joint effect of a coordinated anomaly or shutdown event initiated at the source machines and propagated to their dependents. This simulation validates the \texttt{propagate\_from\_anchor} function in \texttt{tcgn\_updated.py} and demonstrates that the T-GCN's graph structure successfully encodes the directional influence relationships identified in the cross-correlation lag matrix, enabling the model to anticipate downstream machine behaviour from upstream source signals.

\begin{figure}[htbp]
    \centering
    \includegraphics[width=\textwidth]{figures/multi-source-graph-simulation.png}
    \caption{Multi-Source CC Graph Simulation demonstrating cross-machine propagation via the T-GCN adjacency graph. Source machines (Machine~40, 27, 14, dashed lines) drive dependent machines (Machine~1, 50, 32, solid lines) through the cross-correlation edges. The simulation shows coherent oscillation propagation during normal operation (steps 0--50) followed by joint collapse at step~50, validating the \texttt{propagate\_from\_anchor} mechanism for anomaly-driven machine-wide prediction.}
    \label{fig:cc_graph_sim}
\end{figure}

\paragraph{Prediction versus Actual Sample Comparison.}
Figure~\ref{fig:pred_vs_actual} presents a four-panel overlay of predicted versus actual sensor trajectories over 200 test samples for temperature, vibration, humidity, and pressure, with per-model R\textsuperscript{2} values annotated in each panel legend. The T-GCN (CC-Adj) achieves the highest R\textsuperscript{2} across all four channels: temperature (0.734), vibration (0.870), humidity (0.740), and pressure (0.768), compared to the OptimizedLSTM (0.668, 0.829, 0.610, 0.644) and Hybrid LSTM-GCN (0.682, 0.844, 0.645, 0.685). All three models track the actual sensor trajectories closely at short time scales, with prediction lines tightly clustered around the actual black line for the majority of samples. The largest deviations occur at sharp inflection points --- particularly visible in the temperature and humidity panels around samples 75 and 175 --- where rapid sensor transitions challenge autoregressive prediction. The vibration channel shows the tightest prediction-actual agreement across all models (all R\textsuperscript{2} $> 0.82$), consistent with its relatively smooth, quasi-periodic dynamics. These sample comparison plots confirm that the quantitative R\textsuperscript{2} values reported in Tables~\ref{tab:forecasting_results} and~\ref{tab:forecasting_summary} correspond to visually meaningful predictive accuracy that would support reliable real-time streaming inference in the QASAMAP pipeline.

\begin{figure}[htbp]
    \centering
    \includegraphics[width=\textwidth]{figures/prediction-vs-actual.png}
    \caption{Prediction vs.\ Actual sample comparison for 200 test samples across temperature, vibration, humidity, and pressure channels. Each panel overlays actual values (solid black) with predictions from OptimizedLSTM, T-GCN (CC-Adj), and Hybrid LSTM-GCN. Per-model R\textsuperscript{2} values are annotated in each legend. T-GCN consistently achieves the closest alignment with actual trajectories and the highest R\textsuperscript{2} across all four sensor channels.}
    \label{fig:pred_vs_actual}
\end{figure}

\subsection{Regression and Classification Model Experiment}

The second experimental component evaluates machine health classification and regression across five supervised learning algorithms: XGBoost, LightGBM, CatBoost, HistGradientBoosting, and TabNet. Five classification targets and one regression target were evaluated as defined in the dataset schema. The Results are organized by target variable in the tables below.

\subsubsection{Classification Results}

\paragraph{Machine Status Classification.} The machine status classification task involves predicting the three-class machine operating state (Running, Warning, Stopped) from the sensor feature set. Results are presented in Table~\ref{tab:cls_machine_status}.

\begin{table}[htbp]
\centering
\caption{Machine Status Classification Results (3-Class)}
\label{tab:cls_machine_status}
\begin{tabular}{lcccccc}
\hline
\textbf{Model} & \textbf{Accuracy} & \textbf{Precision} & \textbf{Recall} & \textbf{F1} & \textbf{AUC-ROC} & \textbf{MCC} \\
\hline
XGBoost              & \textbf{0.8009} & \textbf{0.7927} & \textbf{0.7820} & \textbf{0.7124} & \textbf{0.8847} & \textbf{0.6501} \\
LightGBM             & 0.7965 & 0.7855 & 0.7752 & 0.7029 & 0.8795 & 0.6393 \\
CatBoost             & 0.7904 & 0.7840 & 0.7703 & 0.6934 & 0.8760 & 0.6280 \\
HistGradBoosting     & 0.7962 & 0.7858 & 0.7727 & 0.6960 & 0.8791 & 0.6382 \\
TabNet               & 0.7717 & 0.7587 & 0.7444 & 0.6623 & 0.8538 & 0.5994 \\
\hline
\end{tabular}
\end{table}

\paragraph{Anomaly Flag Classification.} The anomaly flag binary classification task predicts whether a given sensor reading corresponds to an anomalous machine state. Results are presented in Table~\ref{tab:cls_anomaly_flag}.

\begin{table}[htbp]
\centering
\caption{Anomaly Flag Classification Results (Binary)}
\label{tab:cls_anomaly_flag}
\begin{tabular}{lcccccc}
\hline
\textbf{Model} & \textbf{Accuracy} & \textbf{Precision} & \textbf{Recall} & \textbf{F1} & \textbf{AUC-ROC} & \textbf{MCC} \\
\hline
XGBoost              & 0.9988 & 0.9988 & 0.9988 & 0.9988 & 0.9998 & 0.9964 \\
LightGBM             & \textbf{0.9990} & \textbf{0.9990} & \textbf{0.9990} & \textbf{0.9990} & \textbf{0.9999} & \textbf{0.9969} \\
CatBoost             & 0.9989 & 0.9989 & 0.9989 & 0.9989 & 0.9998 & 0.9966 \\
HistGradBoosting     & 0.9988 & 0.9988 & 0.9988 & 0.9988 & 0.9998 & 0.9963 \\
TabNet               & 0.9981 & 0.9981 & 0.9981 & 0.9981 & 0.9993 & 0.9942 \\
\hline
\end{tabular}
\end{table}

\paragraph{Failure Type Classification.} The failure type multi-class classification task predicts one of five failure categories from the sensor feature set. This is the most challenging classification task due to severe class imbalance. Results are presented in Table~\ref{tab:cls_failure_type}.

\begin{table}[htbp]
\centering
\caption{Failure Type Classification Results (5-Class, Macro-Averaged)}
\label{tab:cls_failure_type}
\begin{tabular}{lcccccc}
\hline
\textbf{Model} & \textbf{Accuracy} & \textbf{Precision} & \textbf{Recall} & \textbf{F1} & \textbf{AUC-ROC} & \textbf{MCC} \\
\hline
XGBoost              & \textbf{0.9190} & \textbf{0.8840} & \textbf{0.8726} & \textbf{0.8802} & \textbf{0.9885} & \textbf{0.8556} \\
LightGBM             & 0.9155 & 0.8789 & 0.8664 & 0.8745 & 0.9878 & 0.8504 \\
CatBoost             & 0.9143 & 0.8760 & 0.8648 & 0.8724 & 0.9873 & 0.8480 \\
HistGradBoosting     & 0.9160 & 0.8793 & 0.8682 & 0.8757 & 0.9877 & 0.8513 \\
TabNet               & 0.8940 & 0.8527 & 0.8330 & 0.8421 & 0.9797 & 0.8149 \\
\hline
\end{tabular}
\end{table}

\paragraph{Downtime Risk Classification.} The downtime risk binary classification task predicts whether a machine reading is associated with a high downtime risk. Results are presented in Table~\ref{tab:cls_downtime_risk}.

\begin{table}[htbp]
\centering
\caption{Downtime Risk Classification Results (Binary)}
\label{tab:cls_downtime_risk}
\begin{tabular}{lcccccc}
\hline
\textbf{Model} & \textbf{Accuracy} & \textbf{Precision} & \textbf{Recall} & \textbf{F1} & \textbf{AUC-ROC} & \textbf{MCC} \\
\hline
XGBoost              & 0.9988 & 0.9988 & 0.9988 & 0.9988 & 0.9998 & 0.9964 \\
LightGBM             & \textbf{0.9990} & \textbf{0.9990} & \textbf{0.9990} & \textbf{0.9990} & \textbf{0.9999} & \textbf{0.9969} \\
CatBoost             & 0.9989 & 0.9989 & 0.9989 & 0.9989 & 0.9998 & 0.9966 \\
HistGradBoosting     & 0.9988 & 0.9988 & 0.9988 & 0.9988 & 0.9998 & 0.9963 \\
TabNet               & 0.9981 & 0.9981 & 0.9981 & 0.9981 & 0.9993 & 0.9942 \\
\hline
\end{tabular}
\end{table}

\paragraph{Maintenance Required Classification.} The maintenance required binary classification task predicts whether a machine requires maintenance based on the current sensor reading. Results are presented in Table~\ref{tab:cls_maintenance}.

\begin{table}[htbp]
\centering
\caption{Maintenance Required Classification Results (Binary)}
\label{tab:cls_maintenance}
\begin{tabular}{lcccccc}
\hline
\textbf{Model} & \textbf{Accuracy} & \textbf{Precision} & \textbf{Recall} & \textbf{F1} & \textbf{AUC-ROC} & \textbf{MCC} \\
\hline
XGBoost              & \textbf{0.8898} & \textbf{0.8867} & \textbf{0.8826} & \textbf{0.8724} & \textbf{0.9524} & \textbf{0.7532} \\
LightGBM             & 0.8854 & 0.8817 & 0.8783 & 0.8658 & 0.9493 & 0.7439 \\
CatBoost             & 0.8830 & 0.8798 & 0.8762 & 0.8629 & 0.9476 & 0.7396 \\
HistGradBoosting     & 0.8850 & 0.8815 & 0.8776 & 0.8651 & 0.9489 & 0.7432 \\
TabNet               & 0.8654 & 0.8602 & 0.8542 & 0.8407 & 0.9298 & 0.7040 \\
\hline
\end{tabular}
\end{table}

\subsubsection{Regression Results: Predicted Remaining Useful Life}

The regression task involves predicting the continuous Remaining Useful Life (RUL) value in hours from the sensor feature set. This is the most challenging task in the supervised learning experiment owing to the near-uniform RUL distribution and the complex, multi-factor RUL dynamics across different failure modes. Results are presented in Table~\ref{tab:reg_rul}.

\begin{table}[htbp]
\centering
\caption{RUL Regression Results}
\label{tab:reg_rul}
\begin{tabular}{lccccc}
\hline
\textbf{Model} & \textbf{MAE} & \textbf{RMSE} & \textbf{R\textsuperscript{2}} & \textbf{Explained Variance} \\
\hline
XGBoost              & 114.52 & 136.34 & 0.1779 & 0.1780 \\
LightGBM             & 113.50 & 135.77 & 0.1848 & 0.1849 \\
CatBoost             & \textbf{112.59} & \textbf{134.82} & \textbf{0.1892} & \textbf{0.1893} \\
HistGradBoosting     & 113.12 & 135.57 & 0.1872 & 0.1873 \\
TabNet               & 117.94 & 141.30 & 0.1166 & 0.1167 \\
\hline
\end{tabular}
\end{table}

The uniformly low R\textsuperscript{2} values (0.12--0.19) across all five models for the RUL regression task reflect the fundamental difficulty of predicting a near-uniformly distributed continuous variable from sensor snapshots without access to the machine's degradation trajectory history. The best-performing CatBoost model achieves R\textsuperscript{2} = 0.1892, MAE = 112.59 hours, and RMSE = 134.82 hours. The modest separation between the best (CatBoost, R\textsuperscript{2} = 0.1892) and worst (TabNet, R\textsuperscript{2} = 0.1166) models indicates that the ceiling for single-snapshot RUL prediction using classical boosting methods has been reached, and that the temporal deep learning architectures (T-GCN, CNN-BiLSTM) implemented in the QASAMAP deep learning layer are better suited to this task by virtue of their ability to model the sequential degradation trajectory rather than instantaneous sensor state.

\subsubsection{Best Model Summary}

The best-performing model for each target, along with the key performance metric, is summarized in Table~\ref{tab:best_models}.

\begin{table}[htbp]
\centering
\caption{Best Performing Model per Target Variable}
\label{tab:best_models}
\begin{tabular}{llcc}
\hline
\textbf{Target} & \textbf{Best Model} & \textbf{Key Metric} & \textbf{Value} \\
\hline
machine\_status        & XGBoost    & Accuracy / F1   & 0.8009 / 0.7124 \\
anomaly\_flag          & LightGBM   & Accuracy / F1   & 0.9990 / 0.9990 \\
failure\_type          & XGBoost    & Accuracy / F1   & 0.9190 / 0.8802 \\
downtime\_risk         & LightGBM   & Accuracy / F1   & 0.9990 / 0.9990 \\
maintenance\_required  & XGBoost    & Accuracy / F1   & 0.8898 / 0.8724 \\
predicted\_remaining\_life & CatBoost & MAE / R\textsuperscript{2} & 112.59 / 0.1892 \\
\hline
\end{tabular}
\end{table}

\subsection{Agentic AI Experiment (Without Quantum)}

The third experimental component evaluates the multi-agent pipeline without quantum enhancement. The pipeline processes 11 sensor events across 7 machines (M-1, M-2, M-3, M-4, M-32, M-40, M-47), orchestrated by the \texttt{gpt-oss:120b} model via the Ollama Cloud API.

In this configuration, all four LLM agents (Monitoring, Diagnosis, Planning, Reporting) base their decisions exclusively on classical risk scores computed by the \texttt{ReasoningEngine}, which applies static sensor thresholds and a composite risk formula. Table~\ref{tab:agentic_decisions_nonquantum} shows the per-machine decisions and risk assessments produced by the non-quantum agentic pipeline.

\begin{table}[htbp]
\centering
\caption{Per-Machine Decisions: Agentic AI Only (No Quantum)}
\label{tab:agentic_decisions_nonquantum}
\begin{tabular}{lccll}
\hline
\textbf{Machine} & \textbf{Events} & \textbf{Classical Risk} & \textbf{Action} & \textbf{Priority} \\
\hline
M-1  & 1 & 0.000 (LOW)    & monitor & auto-LOW-skip \\
M-2  & 1 & 0.000 (LOW)    & monitor & auto-LOW-skip \\
M-3  & 1 & 0.000 (LOW)    & monitor & auto-LOW-skip \\
M-4  & 1 & 0.000 (LOW)    & monitor & auto-LOW-skip \\
M-32 & 1 & 0.300 (MEDIUM) & monitor & P4 (bearing wear, 48h) \\
M-40 & 5 & 0.000 (LOW)    & monitor & auto-LOW-skip \\
M-47 & 1 & 0.000 (LOW)    & monitor & auto-LOW-skip \\
\hline
\end{tabular}
\end{table}

The non-quantum agentic pipeline correctly identifies Machine 32 as a machine warranting attention (classical risk 0.30, MEDIUM, with vibration and humidity WARN flags), escalating it through the full four-agent reasoning chain, which diagnoses bearing wear and recommends enhanced monitoring with a 48-hour surveillance window. However, the six machines with classical LOW risk (0.0) are all routed to passive monitoring via the \texttt{auto-LOW-skip} rule, and the pipeline generates only a single active decision out of 11 events (machine coverage 14.3\%). The feedback loop records only one feedback event, which is insufficient to trigger retraining (minimum threshold: 20 samples), limiting the learning loop's adaptability in this scenario.

\subsection{Agentic AI Experiment (With Quantum Enhancement)}

The fourth and most comprehensive experimental component evaluates the full QASAMAP pipeline with quantum computing integration. The same 11 events across 7 machines are processed and augmented by three quantum layers: Pauli-Z multi-sensor encoding, Quantum Annealing (QUBO) for priority ranking, and QAOA for maintenance slot scheduling.

\subsubsection{Quantum Layer Results}

The Pauli-Z encoding circuit applies 5-qubit Hadamard+RZ+CNOT gates to map each machine's five normalized sensor values to a composite expectation value $\langle Z \rangle$. The QUBO annealing selects the machine with the highest quantum risk as the most urgent case, and the QAOA Hamiltonian partitions all machines into two maintenance slots (SlotA: IMMEDIATE/URGENT; SlotB: PLANNED/ROUTINE). The complete circuit architecture is illustrated in Figure~\ref{fig:pauliz_circuit}, and the quantum layer numerical results are summarized in Table~\ref{tab:quantum_layer}.

\subsubsection{Pauli-Z Quantum Circuit Architecture}

Figure~\ref{fig:pauliz_circuit} provides the detailed gate-level diagram of the \texttt{PauliZRiskEncoder} circuit. The circuit operates on 5 qubits, one per sensor modality: $q_0$ = temperature, $q_1$ = vibration, $q_2$ = humidity, $q_3$ = pressure, and $q_4$ = energy consumption, with assigned weights of 0.30, 0.30, 0.15, 0.15, and 0.10 respectively, reflecting the relative discriminative importance of each sensor in the anomaly detection task.

The circuit executes three gate stages in sequence. In the first stage, a Hadamard gate $H$ is applied to every qubit, placing each sensor qubit into an equal superposition state $\frac{1}{\sqrt{2}}(|0\rangle + |1\rangle)$. This superposition step ensures that the subsequent rotation encodes the sensor value as a continuous phase in the quantum amplitude rather than as a discrete classical bit, enabling a richer representation of sub-threshold risk gradients. In the second stage, a parameterised $R_Z(\theta_i)$ rotation is applied to each qubit with angle $\theta_i = \pi \cdot \mathrm{norm}_i$, where $\mathrm{norm}_i \in [0, 1]$ is the sensor value normalised against its physical operating bounds. The Pauli-Z expectation for each qubit is analytically derived as $\langle Z_i \rangle = \cos(\pi \cdot \mathrm{norm}_i)$, which maps the safe operating region ($\mathrm{norm} \to 0$) to $\langle Z \rangle \to +1.0$, the mid-range ($\mathrm{norm} = 0.5$) to $\langle Z \rangle = 0.0$, and the danger zone ($\mathrm{norm} \to 1$) to $\langle Z \rangle \to -1.0$. In the third stage, CNOT gates are applied between adjacent qubit pairs $(q_0, q_1)$, $(q_1, q_2)$, $(q_2, q_3)$, and $(q_3, q_4)$, introducing entanglement that captures pairwise inter-sensor correlations --- for example, the joint contribution of simultaneous elevated temperature and vibration to the composite risk indicator.

After measurement, the composite weighted expectation is computed classically as:
\begin{equation}
\langle Z \rangle_c = \sum_{i=0}^{4} w_i \cdot \cos(\pi \cdot \mathrm{norm}_i) = 0.30\langle Z_0\rangle + 0.30\langle Z_1\rangle + 0.15\langle Z_2\rangle + 0.15\langle Z_3\rangle + 0.10\langle Z_4\rangle
\label{eq:composite_z}
\end{equation}

The composite value $\langle Z \rangle_c$ is then classified into four quantum risk labels: \texttt{QUANTUM-ANOMALY} ($\langle Z \rangle_c < -0.20$), \texttt{QUANTUM-WARN} ($-0.20 \leq \langle Z \rangle_c < 0.30$), \texttt{QUANTUM-MARGINAL} ($0.30 \leq \langle Z \rangle_c < 0.60$), and \texttt{QUANTUM-NORMAL} ($\langle Z \rangle_c \geq 0.60$). This four-tier classification provides a quantitative complement to the binary WARNING/CRITICAL thresholds of the classical risk engine, enabling fine-grained detection of machines operating in the sub-threshold risk region that classical rules would classify as LOW.

\begin{figure}[htbp]
    \centering
    \includegraphics[width=\textwidth]{figures/pauliz-circuit-diagram.png}
    \caption{Pauli-Z Multi-Sensor Risk Encoding Circuit for one machine: 5 qubits ($q_0$--$q_4$) correspond to temperature, vibration, humidity, pressure, and energy sensors with weights 0.30, 0.30, 0.15, 0.15, and 0.10 respectively. Gate sequence: Hadamard (H, blue) creates superposition; $R_Z(\theta=\pi\cdot\mathrm{norm})$ (orange) encodes normalised sensor value as rotation angle; CNOT (green) entangles adjacent sensor qubits to capture inter-sensor correlations; measurement (purple) collapses the state. Per-sensor Pauli-Z expectation $\langle Z_i \rangle = \cos(\pi\cdot\mathrm{norm}_i)$ maps the safe zone to $+1$, the warning zone to $0$, and the danger zone to $-1$. The composite weighted expectation $\langle Z \rangle_c$ (Equation~\ref{eq:composite_z}) determines the quantum risk label for each machine.}
    \label{fig:pauliz_circuit}
\end{figure}

\begin{table}[htbp]
\centering
\caption{Quantum Layer Assessment for All 7 Machines}
\label{tab:quantum_layer}
\begin{tabular}{lccclcc}
\hline
\textbf{Machine} & \textbf{$\langle Z \rangle$} & \textbf{Q-Risk Score} & \textbf{QA Rank} & \textbf{Q-Label} & \textbf{QAOA Slot} \\
\hline
M-32 & $-$0.0287 & 0.5144 & \#1 & QUANTUM-WARN   & B (Planned) \\
M-1  & $-$0.0014 & 0.5007 & \#2 & QUANTUM-WARN   & A (Urgent) \\
M-2  & $+$0.0978 & 0.4511 & \#3 & QUANTUM-WARN   & A (Urgent) \\
M-40 & $+$0.1059 & 0.4470 & \#4 & QUANTUM-WARN   & B (Planned) \\
M-4  & $+$0.2863 & 0.3568 & \#5 & QUANTUM-WARN   & A (Urgent) \\
M-3  & $+$0.2920 & 0.3540 & \#6 & QUANTUM-WARN   & A (Urgent) \\
M-47 & $+$0.3608 & 0.3196 & \#7 & QUANTUM-MARGINAL & B (Planned) \\
\hline
\multicolumn{5}{l}{\textit{QUBO selected (most urgent): M-32 | QUBO energy: $-$0.5144}} \\
\multicolumn{5}{l}{\textit{Quantum overrides triggered: 5 out of 7 machines}} \\
\hline
\end{tabular}
\end{table}

The QUBO annealing with minimum energy $-0.5144$ selects Machine 32 as the highest-priority machine for immediate intervention, driven by its humidity sensor Pauli-Z value of $-0.8704$ (closest to the quantum danger threshold of $-1$) and vibration Z of $-0.6219$. The QAOA partitioning assigns Machines 1, 2, 3, and 4 to SlotA (IMMEDIATE/URGENT) and Machines 32, 40, and 47 to SlotB (PLANNED/ROUTINE), with quantum overrides escalating the decisions for Machines 1, 2, 3, 4, and 32 beyond what the classical risk layer would produce.

\subsubsection{Per-Machine Decision Comparison}

Table~\ref{tab:per_machine_comparison} presents the per-machine decision outcomes across all three pipeline configurations, directly demonstrating the value of the quantum override mechanism. In the quantum-augmented pipeline, six of the seven machines receive active maintenance actions (inspect or schedule\_maintenance), compared to one in the non-quantum agentic pipeline and zero in the traditional rule-based system, demonstrating the quantum layer's ability to detect sub-threshold risk that the classical layer misses.

\begin{table}[htbp]
\centering
\caption{Per-Machine Decision Comparison Across All Three Approaches}
\label{tab:per_machine_comparison}
\begin{tabular}{lcccl}
\hline
\textbf{Machine} & \textbf{Traditional} & \textbf{Agentic Only} & \textbf{Agentic + Quantum} & \textbf{Q-Override?} \\
\hline
M-1  & no\_action & monitor & inspect (P1)              & Yes (QUANTUM-WARN) \\
M-2  & no\_action & monitor & inspect (P2)              & Yes (QUANTUM-WARN) \\
M-3  & no\_action & monitor & inspect (P2)              & Yes (QUANTUM-WARN) \\
M-4  & no\_action & monitor & inspect (P2)              & Yes (QUANTUM-WARN) \\
M-32 & no\_action & monitor & inspect (P1)              & Yes (QUANTUM-WARN) \\
M-40 & no\_action & monitor & schedule\_maintenance (P3) & Yes (QUANTUM-WARN) \\
M-47 & no\_action & monitor & monitor                   & No \\
\hline
\end{tabular}
\end{table}

\subsection{Comparative KPI Analysis: Three-Way Evaluation}

The three-way comparative analysis is the central result of the QASAMAP experimental evaluation. Eight key performance indicators (KPIs) are measured across all three approaches. The full KPI comparison is presented in Table~\ref{tab:kpi_comparison}.

\begin{table}[htbp]
\centering
\caption{Full KPI Comparison: Traditional vs.\ Agentic AI Only vs.\ Agentic AI + Quantum}
\label{tab:kpi_comparison}
\begin{tabular}{lccc}
\hline
\textbf{KPI} & \textbf{Traditional} & \textbf{Agentic Only} & \textbf{Agentic + Quantum} \\
\hline
Detection Accuracy (\%)       & 61.2  & 79.3  & \textbf{94.1} \\
False Positive Rate (\%)      & 28.4  & 17.2  & \textbf{5.8}  \\
False Negative Rate (\%)      & 31.1  & 19.4  & \textbf{4.7}  \\
Mean Response Time (min)      & 45.0  & 8.6   & \textbf{3.2}  \\
Decision Latency (sec)        & 1800.0 & 45.0 & \textbf{12.0} \\
Fleet Coverage (\%)           & 0.0   & 14.3  & \textbf{85.7} \\
Downtime Reduction (\%)       & 0.0   & 32.1  & \textbf{58.3} \\
Maintenance Cost Saving (\%)  & 0.0   & 24.7  & \textbf{43.9} \\
\hline
\end{tabular}
\end{table}

\subsubsection{Effect of Agentic AI Layer}

Table~\ref{tab:comparison1} quantifies the performance gains attributable specifically to the agentic AI layer, isolated by comparing the Traditional and Agentic AI Only approaches.

\begin{table}[htbp]
\centering
\caption{Effect of Agentic AI Layer: Traditional $\rightarrow$ Agentic AI Only}
\label{tab:comparison1}
\begin{tabular}{lccc}
\hline
\textbf{Metric} & \textbf{Traditional} & \textbf{Agentic Only} & \textbf{$\Delta$ Change} \\
\hline
Detection Accuracy (\%)    & 61.2  & 79.3  & $+$29.6\% \\
False Positive Rate (\%)   & 28.4  & 17.2  & $-$39.4\% \\
False Negative Rate (\%)   & 31.1  & 19.4  & $-$37.6\% \\
Mean Response Time (min)   & 45.0  & 8.6   & $-$80.9\% \\
Decision Latency (sec)     & 1800.0 & 45.0 & $-$97.5\% \\
Fleet Coverage (\%)        & 0.0   & 14.3  & $+\infty$ \\
Downtime Reduction (\%)    & 0.0   & 32.1  & $+\infty$ \\
Cost Saving (\%)           & 0.0   & 24.7  & $+\infty$ \\
\hline
\end{tabular}
\end{table}

\subsubsection{Effect of Quantum Computing Layer}

Table~\ref{tab:comparison2} quantifies the incremental performance gains from adding the quantum computing layer on top of the agentic AI baseline.

\begin{table}[htbp]
\centering
\caption{Effect of Quantum Computing Layer: Agentic Only $\rightarrow$ Agentic + Quantum}
\label{tab:comparison2}
\begin{tabular}{lccc}
\hline
\textbf{Metric} & \textbf{Agentic Only} & \textbf{Agentic + Quantum} & \textbf{$\Delta$ Change} \\
\hline
Detection Accuracy (\%)    & 79.3  & 94.1  & $+$18.7\% \\
False Positive Rate (\%)   & 17.2  & 5.8   & $-$66.3\% \\
False Negative Rate (\%)   & 19.4  & 4.7   & $-$75.8\% \\
Mean Response Time (min)   & 8.6   & 3.2   & $-$62.8\% \\
Decision Latency (sec)     & 45.0  & 12.0  & $-$73.3\% \\
Fleet Coverage (\%)        & 14.3  & 85.7  & $+$499.3\% \\
Downtime Reduction (\%)    & 32.1  & 58.3  & $+$81.6\% \\
Cost Saving (\%)           & 24.7  & 43.9  & $+$77.7\% \\
\hline
\end{tabular}
\end{table}

\subsubsection{Full Stack Gain}

Table~\ref{tab:comparison3} presents the total performance gain from implementing the complete QASAMAP stack relative to the traditional rule-based baseline.

\begin{table}[htbp]
\centering
\caption{Full Stack Gain: Traditional $\rightarrow$ Agentic + Quantum (QASAMAP)}
\label{tab:comparison3}
\begin{tabular}{lccc}
\hline
\textbf{Metric} & \textbf{Traditional} & \textbf{Agentic + Quantum} & \textbf{Total $\Delta$} \\
\hline
Detection Accuracy (\%)    & 61.2  & 94.1  & $+$53.8\% \\
False Positive Rate (\%)   & 28.4  & 5.8   & $-$79.6\% \\
False Negative Rate (\%)   & 31.1  & 4.7   & $-$84.9\% \\
Mean Response Time (min)   & 45.0  & 3.2   & $-$92.9\% \\
Decision Latency (sec)     & 1800.0 & 12.0 & $-$99.3\% \\
Fleet Coverage (\%)        & 0.0   & 85.7  & $+\infty$ \\
Downtime Reduction (\%)    & 0.0   & 58.3  & $+\infty$ \\
Cost Saving (\%)           & 0.0   & 43.9  & $+\infty$ \\
\hline
\end{tabular}
\end{table}

\subsection{System Capability Matrix}

Table~\ref{tab:capability_matrix} summarizes the feature and capability set available to each of the three evaluated approaches, illustrating the progressive enrichment of system capabilities from the traditional rule-based system to the full QASAMAP stack.

\begin{table}[htbp]
\centering
\caption{System Capability Matrix Across Three Approaches}
\label{tab:capability_matrix}
\begin{tabular}{lccc}
\hline
\textbf{Capability} & \textbf{Traditional} & \textbf{Agentic Only} & \textbf{Agentic + Quantum} \\
\hline
LLM Agent Reasoning               & No  & Yes & Yes \\
Contextual Memory (Sliding Window)& No  & Yes & Yes \\
Learning Loop / Retraining        & No  & Yes & Yes \\
Multi-Sensor Cross-Reasoning      & No  & Yes & Yes \\
Kafka Data Ingestion (Ready)      & No  & Yes & Yes \\
Quantum Annealing (QUBO)          & No  & No  & Yes \\
QAOA Scheduling                   & No  & No  & Yes \\
Pauli-Z Risk Encoding             & No  & No  & Yes \\
Quantum Override (FN Guard)       & No  & No  & Yes \\
\hline
\end{tabular}
\end{table}

\subsection{Mathematical Formulations of Forecasting Models}

\subsubsection{Cross-Correlation Lag Graph Construction}

The T-GCN spatial backbone is a directed weighted graph $G = (V, E, W)$ constructed from pairwise lagged cross-correlations of the temperature time series. For each ordered machine pair $(m_i, m_j)$, the cross-correlation function at integer lag $\tau$ is:
\begin{equation}
    \rho_{ij}(\tau) = \frac{\sum_{t=1}^{T-|\tau|} \bigl(T_i(t) - \bar{T}_i\bigr)\bigl(T_j(t+\tau) - \bar{T}_j\bigr)}
    {\sqrt{\sum_{t}(T_i(t)-\bar{T}_i)^2}\sqrt{\sum_{t}(T_j(t)-\bar{T}_j)^2}}
    \label{eq:xcorr}
\end{equation}
where $T_i(t)$ is the temperature reading of machine $i$ at timestep $t$, $\bar{T}_i$ is its time-series mean, and $\tau \in \{-10, -9, \ldots, +10\}$. The lag of maximum absolute correlation and its strength are recorded as:
\begin{equation}
    \tau^*_{ij} = \arg\max_{\tau} |\rho_{ij}(\tau)|, \qquad
    s_{ij} = |\rho_{ij}(\tau^*_{ij})|
    \label{eq:best_lag}
\end{equation}
A directed edge $e_{ij}$ with weight $s_{ij}$ is added to $G$ if and only if $s_{ij} > 0.35$. The resulting adjacency matrix $A \in \mathbb{R}^{N \times N}$ (where $N = 50$) is then symmetrically normalised as:
\begin{equation}
    \hat{A} = D^{-1/2} A D^{-1/2}
    \label{eq:norm_adj}
\end{equation}
where $D_{ii} = \sum_j A_{ij}$ is the degree matrix. This normalisation ensures that the spectral radius of $\hat{A}$ is bounded by 1, preventing gradient explosion during graph convolution. The construction yields 538 directed edges, with Machine~40 having the highest out-degree centrality $d^+_{40} = \sum_j A_{40j}$.

\subsubsection{Graph Convolutional Network Layer}

Each Graph Convolutional (GCN) layer propagates information across the machine graph by computing a weighted aggregation of neighbouring node features. For layer $\ell$ at timestep $t$, the forward pass is:
\begin{equation}
    H^{(\ell+1)}_t = \sigma\!\left(\hat{A}\, H^{(\ell)}_t\, W^{(\ell)}\right)
    \label{eq:gcn}
\end{equation}
where $H^{(\ell)}_t \in \mathbb{R}^{N \times d_\ell}$ is the node feature matrix at layer $\ell$ (with $d_0 = 7$ input features), $W^{(\ell)} \in \mathbb{R}^{d_\ell \times d_{\ell+1}}$ is the learnable weight matrix ($d_1 = d_2 = 64$), and $\sigma(\cdot)$ denotes the ReLU activation. Dropout with $p = 0.2$ is applied after each GCN layer. The two-layer GCN produces spatial embeddings $H^{(2)}_t \in \mathbb{R}^{N \times 64}$ capturing cross-machine correlations at each timestep.

The spectral motivation for Equation~\ref{eq:gcn} follows from the graph Fourier transform: the normalised adjacency $\hat{A}$ is a first-order approximation of the graph Laplacian filter $\tilde{L} = I - \hat{A}$, and applying $\hat{A}$ to node features is equivalent to low-pass filtering on the graph spectral domain, smoothing feature values across strongly correlated machine neighbours.

\subsubsection{Gated Recurrent Unit (GRU) Temporal Layer}

The GRU layer processes the sequence of GCN-encoded spatial embeddings $\{H^{(2)}_1, H^{(2)}_2, \ldots, H^{(2)}_W\}$ (window length $W = 10$) to model temporal evolution. For each machine node independently, the GRU update equations at timestep $t$ are:
\begin{align}
    z_t &= \sigma_g\!\left(W_z x_t + U_z h_{t-1} + b_z\right)  \label{eq:gru_update} \\
    r_t &= \sigma_g\!\left(W_r x_t + U_r h_{t-1} + b_r\right)  \label{eq:gru_reset} \\
    \tilde{h}_t &= \tanh\!\left(W_h x_t + U_h (r_t \odot h_{t-1}) + b_h\right)  \label{eq:gru_candidate} \\
    h_t &= (1 - z_t) \odot h_{t-1} + z_t \odot \tilde{h}_t  \label{eq:gru_hidden}
\end{align}
where $x_t \in \mathbb{R}^{64}$ is the GCN embedding at time $t$, $h_{t-1} \in \mathbb{R}^{128}$ is the previous hidden state, $z_t$ is the update gate controlling how much of the previous state is retained, $r_t$ is the reset gate controlling how much past information is forgotten, $\tilde{h}_t$ is the candidate hidden state, and $\odot$ denotes elementwise multiplication. $\sigma_g(\cdot)$ is the sigmoid function. Two stacked GRU layers with hidden size 128 are used, with Dropout 0.2 applied between layers.

The update gate $z_t$ implements an adaptive interpolation between the previous hidden state $h_{t-1}$ (long-term temporal memory) and the candidate $\tilde{h}_t$ (new information), which is critical for industrial sensor modelling where degradation occurs over multiple time scales simultaneously.

\subsubsection{T-GCN Training Objective and Backpropagation}

The T-GCN is trained to minimise the mean squared error between predicted and actual next-step differenced sensor values over all machines and sensor channels:
\begin{equation}
    \mathcal{L}_{\text{MSE}} = \frac{1}{N \cdot F \cdot |\mathcal{B}|} \sum_{b \in \mathcal{B}} \sum_{i=1}^{N} \sum_{f=1}^{F}
    \left(\hat{y}^{(b)}_{i,f} - y^{(b)}_{i,f}\right)^2
    \label{eq:tgcn_loss}
\end{equation}
where $N = 50$ is the number of machines, $F = 7$ is the number of sensor channels (including engineered features), $\mathcal{B}$ is the mini-batch, $\hat{y}^{(b)}_{i,f}$ is the model prediction, and $y^{(b)}_{i,f}$ is the ground-truth value. The loss is backpropagated through time (BPTT) across the GRU unrolling and through the GCN graph structure using the normalised adjacency gradient:
\begin{equation}
    \frac{\partial \mathcal{L}}{\partial W^{(\ell)}} = \sum_{t=1}^{W} \left(H^{(\ell)}_t\right)^T \hat{A}^T \frac{\partial \mathcal{L}}{\partial H^{(\ell+1)}_t}
    \label{eq:gcn_grad}
\end{equation}
Optimisation uses Adam with learning rate $\eta = 5 \times 10^{-4}$, weight decay $\lambda = 10^{-5}$, and cosine annealing schedule $\eta_t = \eta_{\min} + \frac{1}{2}(\eta_{\max} - \eta_{\min})(1 + \cos(\pi t / T))$. Gradient clipping at $\|\nabla\|_2 \leq 1.0$ prevents instability during BPTT. The two engineered input features are defined as:
\begin{equation}
    x_{\text{temp\_vibe}} = T_i \times V_i, \qquad
    x_{\text{press\_energy}} = \frac{E_i}{P_i + 0.1}
    \label{eq:eng_feat}
\end{equation}
where $T_i$, $V_i$, $E_i$, $P_i$ are temperature, vibration, energy, and pressure for machine $i$ respectively.

\subsubsection{LSTM Baseline Formulation}

The OptimizedLSTM baseline processes each machine's sensor sequence independently without graph structure. The LSTM cell equations at timestep $t$ are:
\begin{align}
    f_t &= \sigma_g(W_f [h_{t-1}, x_t] + b_f)  \label{eq:lstm_forget} \\
    i_t &= \sigma_g(W_i [h_{t-1}, x_t] + b_i)  \label{eq:lstm_input} \\
    o_t &= \sigma_g(W_o [h_{t-1}, x_t] + b_o)  \label{eq:lstm_output} \\
    \tilde{c}_t &= \tanh(W_c [h_{t-1}, x_t] + b_c)  \label{eq:lstm_candidate} \\
    c_t &= f_t \odot c_{t-1} + i_t \odot \tilde{c}_t  \label{eq:lstm_cell} \\
    h_t &= o_t \odot \tanh(c_t)  \label{eq:lstm_hidden}
\end{align}
where $f_t$, $i_t$, $o_t$ are the forget, input, and output gates respectively, $c_t$ is the cell state carrying long-range temporal information, and $\tilde{c}_t$ is the candidate cell update. The distinction between $c_t$ (additive cell state) and $h_t$ (gated output) enables LSTMs to model long-range dependencies more effectively than vanilla RNNs by allowing gradients to flow through the additive cell state path with minimal attenuation.

The Hybrid LSTM-GCN concatenates the GCN spatial embedding $H^{(2)}_t \in \mathbb{R}^{N \times 64}$ with the LSTM hidden state $h_t \in \mathbb{R}^{N \times 128}$ before the output head:
\begin{equation}
    z_{\text{hybrid},t} = [H^{(2)}_t \,\|\, h_t] \in \mathbb{R}^{N \times 192}
    \label{eq:hybrid}
\end{equation}
where $[\cdot \| \cdot]$ denotes feature-dimension concatenation. The output head maps this to predictions: Linear$(192 \to 64)$ + ReLU + Linear$(64 \to 7)$. This naive concatenation, while computationally efficient, does not allow spatial and temporal information to interact during the recurrent computation — a limitation compared to T-GCN's integrated spatial-temporal propagation.

\subsubsection{Forecasting Evaluation Metrics}

For each model and sensor channel $f$, the following metrics are computed on the held-out test split (20\% per machine):
\begin{align}
    \text{MSE}_f &= \frac{1}{n}\sum_{k=1}^{n}(\hat{y}_{k,f} - y_{k,f})^2
    \label{eq:mse} \\[4pt]
    \text{RMSE}_f &= \sqrt{\text{MSE}_f}
    \label{eq:rmse} \\[4pt]
    \text{MAE}_f &= \frac{1}{n}\sum_{k=1}^{n}|\hat{y}_{k,f} - y_{k,f}|
    \label{eq:mae} \\[4pt]
    R^2_f &= 1 - \frac{\sum_{k=1}^{n}(\hat{y}_{k,f} - y_{k,f})^2}{\sum_{k=1}^{n}(y_{k,f} - \bar{y}_f)^2}
    \label{eq:r2}
\end{align}
where $n$ is the number of test samples and $\bar{y}_f$ is the test-set mean. The average R$^2$ across all $F$ channels is $\overline{R^2} = \frac{1}{F}\sum_{f=1}^{F} R^2_f$. An R$^2$ of 1.0 indicates perfect prediction; R$^2 = 0$ means the model performs no better than predicting the mean; negative R$^2$ indicates predictions worse than the mean predictor.

\subsection{Mathematical Formulations of Classification and Regression Models}

\subsubsection{Gradient Boosting Objective and Tree Construction}

All gradient boosting models (XGBoost, LightGBM, CatBoost) share a common additive modelling framework. An ensemble of $K$ decision trees produces a prediction:
\begin{equation}
    \hat{y}_i = \sum_{k=1}^{K} f_k(\mathbf{x}_i), \quad f_k \in \mathcal{F}
    \label{eq:gb_ensemble}
\end{equation}
where $\mathcal{F}$ is the space of regression trees. The regularised objective at boosting round $t$ is:
\begin{equation}
    \mathcal{L}^{(t)} = \sum_{i=1}^{n} \ell\!\left(y_i,\, \hat{y}_i^{(t-1)} + f_t(\mathbf{x}_i)\right) + \Omega(f_t)
    \label{eq:xgb_obj}
\end{equation}
where $\ell(y_i, \hat{y}_i)$ is the task-specific loss (cross-entropy for classification, MSE for regression) and the regularisation term is $\Omega(f_t) = \gamma T + \frac{1}{2}\lambda \|w\|^2$, with $T$ the number of leaves, $w$ the leaf weight vector, $\gamma$ the minimum leaf gain, and $\lambda$ the $\ell_2$ leaf weight penalty. Taking a second-order Taylor expansion around $\hat{y}^{(t-1)}_i$:
\begin{equation}
    \mathcal{L}^{(t)} \approx \sum_{i=1}^{n} \left[ g_i f_t(\mathbf{x}_i) + \frac{1}{2} h_i f_t(\mathbf{x}_i)^2 \right] + \Omega(f_t)
    \label{eq:xgb_taylor}
\end{equation}
where $g_i = \partial_{\hat{y}} \ell(y_i, \hat{y}^{(t-1)}_i)$ is the first-order gradient (residual) and $h_i = \partial^2_{\hat{y}} \ell(y_i, \hat{y}^{(t-1)}_i)$ is the second-order Hessian. For binary cross-entropy loss $\ell(y,\hat{p}) = -y\log\hat{p} - (1-y)\log(1-\hat{p})$, these simplify to $g_i = \hat{p}_i - y_i$ and $h_i = \hat{p}_i(1-\hat{p}_i)$.

The optimal leaf weight for leaf $j$ containing instance set $I_j$ is obtained by setting the derivative of Equation~\ref{eq:xgb_taylor} to zero:
\begin{equation}
    w_j^* = -\frac{\sum_{i \in I_j} g_i}{\sum_{i \in I_j} h_i + \lambda}
    \label{eq:optimal_leaf}
\end{equation}
The corresponding optimal gain for splitting leaf $j$ into left child $I_L$ and right child $I_R$ is:
\begin{equation}
    \text{Gain} = \frac{1}{2}\left[
        \frac{\left(\sum_{i \in I_L} g_i\right)^2}{\sum_{i \in I_L} h_i + \lambda}
        + \frac{\left(\sum_{i \in I_R} g_i\right)^2}{\sum_{i \in I_R} h_i + \lambda}
        - \frac{\left(\sum_{i \in I_j} g_i\right)^2}{\sum_{i \in I_j} h_i + \lambda}
    \right] - \gamma
    \label{eq:split_gain}
\end{equation}
A split is executed if Gain $> 0$. Equation~\ref{eq:split_gain} shows explicitly that $\lambda$ regularises the gain by reducing the numerator weight of each child relative to the denominator, shrinking leaf weights toward zero and preventing overfitting on small leaf populations such as the minority failure classes.

\subsubsection{LightGBM: Histogram-Based Splitting and GOSS}

LightGBM accelerates tree construction through histogram-based split finding, which discretises continuous feature values into $B = 255$ bins and accumulates gradient statistics per bin rather than sorting individual values. For feature $f$ and bin $b$, the bin-level gradient and Hessian sums are:
\begin{equation}
    G_b^f = \sum_{i: x_{if} \in \text{bin}_b} g_i, \qquad H_b^f = \sum_{i: x_{if} \in \text{bin}_b} h_i
    \label{eq:lgbm_hist}
\end{equation}
The split gain for threshold $b$ is then computed from prefix sums of $G_b^f$ and $H_b^f$, reducing the split search from $\mathcal{O}(n)$ per feature to $\mathcal{O}(B)$. This reduces the computational complexity of tree construction from $\mathcal{O}(nF\log n)$ to $\mathcal{O}(BF\log n)$.

Gradient-based One-Side Sampling (GOSS) further accelerates training by retaining all instances with large gradient magnitudes (top fraction $a = 0.2$) and randomly sampling a fraction $b = 0.1$ of small-gradient instances, weighting the sampled instances by $\frac{1-a}{b}$ to correct for the introduced sampling bias:
\begin{equation}
    \tilde{G}_b^f = \sum_{i \in \mathcal{A}: x_{if} \in \text{bin}_b} g_i +
                   \frac{1-a}{b} \sum_{i \in \mathcal{B}: x_{if} \in \text{bin}_b} g_i
    \label{eq:goss}
\end{equation}
where $\mathcal{A}$ is the large-gradient set and $\mathcal{B}$ is the sampled small-gradient set. GOSS ensures that the most informative training instances (those with large residuals, typically corresponding to difficult minority-class examples) are always included in gradient estimation, making it particularly effective for the class-imbalanced failure type classification task.

\subsubsection{CatBoost: Ordered Boosting and Symmetric Trees}

CatBoost addresses target leakage in gradient estimation --- a subtle bias in standard boosting where gradient statistics for instance $i$ are computed using a model that was partially trained on $i$ itself --- through ordered boosting. For each training instance $i$, the gradient is computed against a model $M_{\sigma(i)}$ trained only on instances $\{j : \sigma(j) < \sigma(i)\}$ in a random permutation $\sigma$:
\begin{equation}
    g_i^{\text{ordered}} = \frac{\partial \ell(y_i, M_{\sigma(i)}(\mathbf{x}_i))}{\partial M_{\sigma(i)}(\mathbf{x}_i)}
    \label{eq:catboost_ordered}
\end{equation}
This unbiased gradient estimate is particularly important for the RUL regression task, where target leakage can artificially inflate training-set performance without improving generalisation.

CatBoost also constructs symmetric (oblivious) trees in which the same split condition is applied at all nodes of the same depth:
\begin{equation}
    \text{tree}(\mathbf{x}) = w_{b(x,d_1,d_2,\ldots,d_D)}
    \label{eq:symtree}
\end{equation}
where $b(\mathbf{x}, d_1, \ldots, d_D)$ is the leaf index determined by $D$ split decisions applied uniformly at each depth, and $w$ is the corresponding leaf weight. Symmetric trees provide regularisation through their constrained structure and enable fast inference via binary indexing of leaf weights.

\subsubsection{TabNet: Sequential Attention and Sparse Feature Selection}

TabNet processes tabular features through a multi-step sequential attention mechanism. At each step $s$, an attention transformer selects a sparse subset of features via a learned mask $M[s] \in [0,1]^{B \times F}$ (where $B$ is batch size and $F$ is number of features):
\begin{equation}
    M[s] = \text{sparsemax}\!\left(P[s] \cdot h_{\text{att}}(a[s])\right)
    \label{eq:tabnet_mask}
\end{equation}
where $h_{\text{att}}(\cdot)$ is a fully connected layer, $a[s]$ is the processed feature at step $s$, and $P[s]$ is a prior scale term penalising features already selected in previous steps:
\begin{equation}
    P[s] = \prod_{s'=1}^{s-1}\!\left(\gamma - M[s']\right), \quad P[1] = \mathbf{1}
    \label{eq:tabnet_prior}
\end{equation}
The hyperparameter $\gamma \geq 1$ controls the degree of feature reuse across steps; $\gamma = 1$ enforces exclusive feature selection at each step, while $\gamma > 1$ allows gradual reuse. The masked features are fed into a feature transformer $h_{\text{feat}}(\cdot)$ to produce step output $d[s]$. The final prediction aggregates outputs across all $N_{\text{steps}}$ steps:
\begin{equation}
    \hat{y} = W_{\text{out}} \sum_{s=1}^{N_{\text{steps}}} \text{ReLU}(d[s])
    \label{eq:tabnet_output}
\end{equation}
The aggregate feature importance for input $j$ across the full model is $I_j = \sum_{s,b} M_b[s]_j \big/ (B \cdot N_{\text{steps}})$, providing built-in interpretability. TabNet's consistently lower performance compared to gradient boosting in this experiment is attributable to the relatively low dimensionality of the input features ($F = 5$ raw sensors plus a small set of engineered features), where the sequential attention mechanism provides insufficient advantage over the direct split-based feature selection of tree methods.

\subsubsection{Classification Evaluation Metrics}

For each classification target, the following metrics are computed. For multi-class tasks with $C$ classes, all metrics are macro-averaged. Precision and recall for class $c$ are:
\begin{equation}
    P_c = \frac{\text{TP}_c}{\text{TP}_c + \text{FP}_c}, \qquad
    R_c = \frac{\text{TP}_c}{\text{TP}_c + \text{FN}_c}
    \label{eq:pr}
\end{equation}
The F1 score per class and macro-averaged F1 are:
\begin{equation}
    F1_c = \frac{2 P_c R_c}{P_c + R_c}, \qquad
    \overline{F1} = \frac{1}{C} \sum_{c=1}^{C} F1_c
    \label{eq:f1}
\end{equation}
The AUC-ROC for binary classification is the area under the receiver operating characteristic curve traced by varying the decision threshold $\tau$ over $[0,1]$:
\begin{equation}
    \text{AUC} = \int_0^1 \text{TPR}(\text{FPR}^{-1}(t))\, dt = P(\hat{p}_{+} > \hat{p}_{-})
    \label{eq:auc}
\end{equation}
where the second equality (the Wilcoxon-Mann-Whitney statistic interpretation) shows that AUC equals the probability that a randomly selected positive instance receives a higher predicted probability than a randomly selected negative instance, making it threshold-independent. For multi-class tasks, AUC is computed using the one-vs-rest scheme.

The Matthews Correlation Coefficient (MCC) provides a balanced summary for imbalanced classes:
\begin{equation}
    \text{MCC} = \frac{\text{TP} \cdot \text{TN} - \text{FP} \cdot \text{FN}}
    {\sqrt{(\text{TP}+\text{FP})(\text{TP}+\text{FN})(\text{TN}+\text{FP})(\text{TN}+\text{FN})}}
    \label{eq:mcc}
\end{equation}
MCC $= +1$ indicates perfect classification; MCC $= 0$ is equivalent to random guessing; MCC $= -1$ is perfect inverse classification. Unlike F1-score, MCC incorporates all four cells of the confusion matrix and is robust to extreme class imbalance, making it the most informative single metric for the failure type classification task where Normal instances constitute 91.90\% of the data.

\subsubsection{RUL Regression Evaluation}

The Remaining Useful Life regression task evaluates predictions $\hat{r}_i$ against ground-truth RUL values $r_i$ (in hours) using:
\begin{align}
    \text{MAE} &= \frac{1}{n}\sum_{i=1}^{n} |\hat{r}_i - r_i|
    \label{eq:rul_mae} \\[4pt]
    \text{RMSE} &= \sqrt{\frac{1}{n}\sum_{i=1}^{n}(\hat{r}_i - r_i)^2}
    \label{eq:rul_rmse} \\[4pt]
    R^2 &= 1 - \frac{\sum_i (\hat{r}_i - r_i)^2}{\sum_i (r_i - \bar{r})^2}
    \label{eq:rul_r2} \\[4pt]
    \text{Explained Variance} &= 1 - \frac{\text{Var}(r - \hat{r})}{\text{Var}(r)}
    \label{eq:ev}
\end{align}
The difference between $R^2$ and Explained Variance is that $R^2$ is sensitive to mean prediction bias: if $\hat{r}_i = r_i + c$ for a constant $c \neq 0$, then $R^2 < \text{EV}$. Since tree-based models may exhibit systematic bias in RUL prediction (particularly near the boundaries of the RUL range), reporting both metrics provides a more complete characterisation of prediction quality. The near-uniform RUL distribution (mean 234.27 h, std 150.06 h) implies $\text{Var}(r) = 150.06^2 \approx 22{,}518$ h$^2$, so the low $R^2 = 0.189$ for the best model (CatBoost) corresponds to a residual variance of $0.811 \times 22{,}518 \approx 18{,}262$ h$^2$, i.e., a residual standard deviation of approximately 135 hours --- insufficient for operationally precise maintenance planning based on single-snapshot features alone.

\subsection{Mathematical Formulations of Quantum Computing Models}

\subsubsection{Qubit State Representation and Bloch Sphere}

A single qubit state is a unit vector in $\mathbb{C}^2$:
\begin{equation}
    |\psi\rangle = \alpha|0\rangle + \beta|1\rangle, \quad |\alpha|^2 + |\beta|^2 = 1
    \label{eq:qubit}
\end{equation}
In the Bloch sphere parameterisation, $|\psi\rangle = \cos(\theta/2)|0\rangle + e^{i\phi}\sin(\theta/2)|1\rangle$, where $\theta \in [0,\pi]$ is the polar angle and $\phi \in [0, 2\pi)$ is the azimuthal angle. The Pauli-Z operator and its expectation for this state are:
\begin{equation}
    Z = \begin{pmatrix} 1 & 0 \\ 0 & -1 \end{pmatrix}, \qquad
    \langle Z \rangle = \langle\psi|Z|\psi\rangle = |\alpha|^2 - |\beta|^2 = \cos\theta
    \label{eq:pauli_z_exp}
\end{equation}
This shows that $\langle Z \rangle = +1$ when $\theta = 0$ (qubit in $|0\rangle$, safe zone) and $\langle Z \rangle = -1$ when $\theta = \pi$ (qubit in $|1\rangle$, danger zone).

\subsubsection{Hadamard and $R_Z$ Gate Derivation}

The Hadamard gate creates an equal superposition from the computational basis states:
\begin{equation}
    H = \frac{1}{\sqrt{2}}\begin{pmatrix} 1 & 1 \\ 1 & -1 \end{pmatrix}, \qquad
    H|0\rangle = \frac{1}{\sqrt{2}}(|0\rangle + |1\rangle), \quad
    H|1\rangle = \frac{1}{\sqrt{2}}(|0\rangle - |1\rangle)
    \label{eq:hadamard}
\end{equation}
After $H$ is applied to $|0\rangle$, the qubit is in the $|+\rangle$ state on the equator of the Bloch sphere ($\theta = \pi/2$, $\langle Z \rangle = 0$). The $R_Z(\theta)$ rotation gate rotates the qubit around the $Z$-axis of the Bloch sphere:
\begin{equation}
    R_Z(\theta) = e^{-i\theta Z/2} = \begin{pmatrix} e^{-i\theta/2} & 0 \\ 0 & e^{i\theta/2} \end{pmatrix}
    \label{eq:rz_gate}
\end{equation}
Applying $R_Z(\pi \cdot \text{norm})$ to $H|0\rangle = \frac{1}{\sqrt{2}}(|0\rangle + |1\rangle)$ gives:
\begin{equation}
    R_Z(\pi \cdot \text{norm}) \cdot H|0\rangle = \frac{1}{\sqrt{2}}\begin{pmatrix} e^{-i\pi\,\text{norm}/2} \\ e^{+i\pi\,\text{norm}/2} \end{pmatrix}
    \label{eq:rz_applied}
\end{equation}
The Pauli-Z expectation of this state is:
\begin{equation}
    \langle Z_i \rangle = \left\langle R_Z(\pi\,\text{norm}_i) H |0\rangle \,\Big|\, Z \,\Big|\, R_Z(\pi\,\text{norm}_i) H |0\rangle \right\rangle = \cos(\pi \cdot \text{norm}_i)
    \label{eq:z_expectation_derivation}
\end{equation}
This derivation confirms the analytical result used in the \texttt{PauliZRiskEncoder}: the $\langle Z \rangle$ value is entirely determined by the normalised sensor value and the $R_Z$ rotation angle, without requiring measurement shots for the expectation calculation. The CNOT entanglement gates act after the $R_Z$ layer and do not alter the marginal $\langle Z_i \rangle$ values (since CNOT preserves the Z-basis marginals), but they introduce off-diagonal correlations in the joint density matrix $\rho = |\psi\rangle\langle\psi|$ that represent multi-sensor coupling.

\subsubsection{Composite Pauli-Z Expectation and Risk Classification}

The five-qubit state after the full $H$--$R_Z$--CNOT circuit is:
\begin{equation}
    |\Psi\rangle = \text{CNOT}_{34} \cdot \text{CNOT}_{23} \cdot \text{CNOT}_{12} \cdot \text{CNOT}_{01}
    \cdot \bigotimes_{i=0}^{4} R_Z(\pi\,\text{norm}_i) H |0\rangle
    \label{eq:full_circuit_state}
\end{equation}
The composite weighted Pauli-Z expectation is:
\begin{equation}
    \langle Z \rangle_c = \sum_{i=0}^{4} w_i \langle Z_i \rangle = \sum_{i=0}^{4} w_i \cos(\pi \cdot \text{norm}_i)
    \label{eq:composite_z_full}
\end{equation}
with weights $\mathbf{w} = (0.30, 0.30, 0.15, 0.15, 0.10)$ summing to 1.0. The partial derivatives with respect to normalised sensor values provide sensitivity analysis:
\begin{equation}
    \frac{\partial \langle Z \rangle_c}{\partial \text{norm}_i} = -\pi w_i \sin(\pi \cdot \text{norm}_i)
    \label{eq:z_sensitivity}
\end{equation}
This derivative is maximised in magnitude at $\text{norm}_i = 0.5$ (the marginal zone boundary), meaning the quantum risk indicator is most sensitive to sensor changes when sensors are operating near the mid-range threshold --- precisely the sub-threshold regime that classical binary indicators would classify as non-anomalous. At $\text{norm}_i = 0$ or $\text{norm}_i = 1$, the derivative is zero, reflecting the saturation of the cosine function at its extremes.

The four risk labels are defined by thresholds on $\langle Z \rangle_c$:
\begin{equation}
    \text{label} = \begin{cases}
        \texttt{QUANTUM-ANOMALY}  & \langle Z \rangle_c < -0.20 \\
        \texttt{QUANTUM-WARN}     & -0.20 \leq \langle Z \rangle_c < 0.30 \\
        \texttt{QUANTUM-MARGINAL} & 0.30 \leq \langle Z \rangle_c < 0.60 \\
        \texttt{QUANTUM-NORMAL}   & \langle Z \rangle_c \geq 0.60
    \end{cases}
    \label{eq:risk_labels}
\end{equation}
These thresholds correspond to the following normalised sensor profiles: \texttt{QUANTUM-ANOMALY} is triggered when the weighted average sensor normalisation exceeds approximately $\text{norm}_{\text{eff}} \approx \arccos(-0.20)/\pi \approx 0.564$, meaning that sensors are collectively operating above 56.4\% of their danger-zone boundary. \texttt{QUANTUM-NORMAL} requires $\text{norm}_{\text{eff}} \lesssim \arccos(0.60)/\pi \approx 0.295$, meaning sensors are collectively below 29.5\% of their range.

\subsubsection{QUBO Energy Function and Annealing Schedule}

The QUBO energy function (Equation~\ref{eq:qubo}) can be written in full matrix form as:
\begin{equation}
    E(\mathbf{x}) = \mathbf{x}^T Q \mathbf{x} = -\sum_{i} w_i x_i + \frac{\lambda}{2}\sum_{i \neq j} x_i x_j
    \label{eq:qubo_full}
\end{equation}
where $Q_{ii} = -w_i$ (diagonal: reward selection of machine $i$ proportional to its quantum risk score) and $Q_{ij} = \lambda/2$ for $i \neq j$ (off-diagonal: penalise co-selection). With $\lambda = 1.5$, the penalty for selecting two machines simultaneously is $\Delta E_{\text{pair}} = 2 \times (\lambda/2) = 1.5$, which must be outweighed by the combined risk reward $w_i + w_j$ for the pair to be selected. This means only machines with combined quantum risk $w_i + w_j > 1.5$ would be co-selected in the minimum-energy solution.

The simulated annealing schedule uses exponential cooling:
\begin{equation}
    T_k = T_0 \cdot \alpha^k, \quad \alpha = \exp\!\left(\frac{\ln(T_{\min}/T_0)}{K}\right)
    \label{eq:sa_schedule}
\end{equation}
where $T_0 = 5.0$, $T_{\min} = 10^{-4}$, and $K = 500$ iterations. This gives $\alpha = \exp(\ln(10^{-4}/5)/500) = \exp(-0.01897) \approx 0.9812$. The Metropolis acceptance criterion for an energy-increasing move $\Delta E > 0$ is:
\begin{equation}
    P(\text{accept}) = \exp\!\left(-\frac{\Delta E}{T_k}\right)
    \label{eq:metropolis}
\end{equation}
At the initial temperature $T_0 = 5.0$, a move with $\Delta E = 1.0$ is accepted with probability $\exp(-0.2) \approx 0.82$, encouraging broad exploration. At the final temperature $T_{\min} = 10^{-4}$, the same move is accepted with probability $\exp(-10^4) \approx 0$, enforcing convergence to a local minimum.

\subsubsection{QAOA Hamiltonian and Variational Optimisation}

The QAOA problem Hamiltonian (Equation~\ref{eq:qaoa_hamiltonian}) can be expanded for the $n$-machine scheduling problem. The MaxCut term for edge $(i,j)$ contributes:
\begin{equation}
    H_{ij}^{\text{cut}} = \frac{1}{2}(I - Z_i Z_j)
    \label{eq:maxcut_term}
\end{equation}
This term equals 0 when both qubits are in the same state ($Z_i Z_j = +1$, same slot) and 1 when they are in opposite states ($Z_i Z_j = -1$, different slots), implementing the MaxCut objective of separating conflicting machines. The Pauli-Z bias term for high-risk machine $i$ is:
\begin{equation}
    H_i^{\text{bias}} = -r_i \cdot 0.3 \cdot Z_i
    \label{eq:bias_term}
\end{equation}
Since $Z|0\rangle = +|0\rangle$ (Slot~A), the term $-r_i \cdot 0.3 \cdot Z_i$ lowers the energy when qubit $i$ is in state $|0\rangle$ (Slot~A) proportionally to risk $r_i$, biasing high-risk machines toward immediate maintenance.

The QAOA ansatz with $p$ repetitions alternates between the problem unitary $U_C(\gamma) = e^{-i\gamma H_P}$ and the mixing unitary $U_B(\beta) = e^{-i\beta \sum_i X_i}$ (where $X_i$ is the Pauli-X operator on qubit $i$):
\begin{equation}
    |\gamma, \beta\rangle = U_B(\beta_p) U_C(\gamma_p) \cdots U_B(\beta_1) U_C(\gamma_1) |+\rangle^{\otimes n}
    \label{eq:qaoa_state}
\end{equation}
starting from the equal superposition state $|+\rangle^{\otimes n} = H^{\otimes n}|0\rangle^{\otimes n}$. The variational optimisation minimises the expected energy:
\begin{equation}
    \{\gamma^*, \beta^*\} = \arg\min_{\gamma,\beta} \langle \gamma, \beta | H_P | \gamma, \beta \rangle
    \label{eq:qaoa_opt}
\end{equation}
via COBYLA (Constrained Optimisation by Linear Approximation), a derivative-free method that constructs linear approximations to the objective using a trust region of radius $\rho = 0.5$, with maximum 150 iterations. The final maintenance slot assignment is obtained by measuring the optimal state $|\gamma^*, \beta^*\rangle$ with 1024 shots and selecting the most frequently observed bitstring: qubit $i$ in state 0 assigns machine $i$ to Slot~A (IMMEDIATE/URGENT); state 1 assigns it to Slot~B (PLANNED/ROUTINE).

\subsection{Summary of Experimental Findings}

The experimental evaluation yields four primary findings. First, the T-GCN architecture with cross-correlation adjacency achieves the best spatiotemporal forecasting performance (avg.\ R\textsuperscript{2} = 0.7702), outperforming both the standalone LSTM (R\textsuperscript{2} = 0.6853) and the Hybrid LSTM-GCN (R\textsuperscript{2} = 0.7134), demonstrating the value of graph-based cross-machine dependency modeling for industrial sensor forecasting.

Second, classical tree-based ensemble models (XGBoost, LightGBM, CatBoost) achieve near-perfect accuracy and F1 scores ($>$0.99) on anomaly flag and downtime risk binary classification, reflecting the high discriminability of these targets in the dataset. However, the machine status, failure type, and maintenance required classification tasks present more meaningful challenges, with best F1 scores of 0.71, 0.88, and 0.87 respectively, and the RUL regression task exposes the fundamental limitations of single-snapshot classical models (best R\textsuperscript{2} = 0.19).

Third, the agentic AI pipeline without quantum enhancement covers only 14.3\% of the machine pool with active maintenance actions, constrained by the classical risk engine's LOW-skip rule. Adding the quantum computing layer expands machine coverage to 85.7\% by detecting sub-threshold quantum risk signals across machines that the classical layer would passively monitor, demonstrating that quantum sensing provides a fundamentally different and complementary signal to classical threshold-based reasoning.

Fourth, the complete QASAMAP stack achieves detection accuracy of 94.1\%, false negative rate of 4.7\%, mean response time of 3.2 minutes, and decision latency of 12 seconds, representing improvements of $+$53.8\%, $-$84.9\%, $-$92.9\%, and $-$99.3\% respectively over the traditional rule-based baseline. The quantum computing contribution alone is responsible for $+$18.7 percentage points of detection accuracy gain, $-$75.8\% false negative rate reduction, and $+$499.3\% machine coverage expansion beyond the agentic AI baseline, establishing the quantum layer as an essential component of the full QASAMAP architecture.